% ════════════════════════════════════════════════════════════════════════════
\section{Supplement to the classification}
\label{app:classify}
% ════════════════════════════════════════════════════════════════════════════

\noindent This appendix carries what \S\ref{sec:classify} summarizes: the
closure condition on the delegated route, the full self-destruct analysis with
the vendor-by-vendor record and the companion decoy, the Obscurity Paradox, and
the registry method --- the coding tables, the jurisdictional comparison with its
confounds, the proxy-victim finding, and the two case analyses.

\subsection{When the remote-delegate escape closes}
\label{app:classify-escape}

\begin{remark}[When the escape closes]\label{rem:escape-closes}
The remedy holds only while the delegate is genuinely remote. If the delegate is
a friend who happens to be visiting during the attack, \emph{both} problems
apply at once --- custody was already non-individual, and the racer is now
co-located and reachable --- so the design becomes deadly-race-susceptible
again. The user is thus forced into an awkward corner: trust someone they can
rarely be near, and accept a perverse incentive to avoid physical proximity to a
trusted friend.
\end{remark}

\subsection{BitVault}
\label{app:classify-bitvault}

\begin{remark}[Sources, and the four columns]\label{rem:bitvault}
\emph{Sources.} No authoritative specification is published. The vendor describes
the product as ``100\% Open Source,'' but its code link resolves to an
organization with no public repositories (archived at
\url{https://archive.is/ONgc4}), and its documentation and audit are served
through a gated document-sharing service. What we have is the vendor's pages, a
concept whitepaper~\cite{bssl} that is self-attributed and hosted outside the
vendor's own links, and correspondence of 26--27 August 2026 (all accessed
1 September 2026); they do not agree in every particular. This bounds what can be
claimed and is \emph{not} an allegation: we have no evidence that any material was
withdrawn, and observed none at any earlier time. The whitepaper's
Taproot/Miniscript policy gates the primary path on \texttt{older()} and both
fallbacks on \texttt{after()}, so the primary delay is a maturation from funding
rather than a window on a request, and the multi-year locks do not re-arm.

\emph{The columns.} Property~\ref{p:know} we leave open: the documentation assumes
hardware and offers an optional ``copy of $A$ kept in a separate location,'' but
nothing in the construction forbids running it from memorized seeds and
re-deriving on demand, which would be deviceless. Property~\ref{p:indiv} fails in
both directions --- the optional off-chain gatekeeper can deny a user the quick
path by declining to release the custodian's co-signature, and past the three-year
lock a custodian and that gatekeeper can spend together. Property~\ref{p:coerc}
fails by the two cases above.

Property~\ref{p:nonob} we also leave open, for a reason worth stating generally:
\textbf{obscurity does not advertise itself}, so classifying a design as
obscurity-dependent is always an inference from what it would have to assume in
order to make sense. Here the deadly-race exposure is plain on the published
policy, plain enough that a Kerckhoffs-aware design would not ship it; so the
design \emph{suggests} an $\Attacker$ assumed ignorant of it --- a reading the
scarcity of public-facing documentation would fit. The only other
reading we can construct is an $\Attacker$ assumed willing to coerce but not to
kill --- which is not a milder premise but a wager on his $c$
(Definition~\ref{def:priced}), staked on the holder's life. Undetermined is the
honest entry between the two, and it costs nothing: the coercion result stands on
structure either way.
\end{remark}

\subsection{The self-destruct class in full}
\label{app:classify-selfdestruct}

\begin{remark}[The obscure self-destruct is the obscure scheme with no winning branch]
\label{rem:two-ignorances}
Obscure defenses are usually compared by \emph{how likely} $\Attacker$ is to be
ignorant. The sharper comparison is \emph{how many} independent ignorances each
one needs, and here the self-destruct is strictly the worse bet --- worse, in
particular, than the decoy it is commonly sold beside.

A decoy needs one premise: that $\Attacker$ does not know decoys exist. Grant him
that ignorance and the decoy \emph{terminates the encounter} --- he leaves with a
plausible balance and a story that closes.

A self-destruct needs two, and they are independent. It needs $\Attacker$ ignorant
of the feature (sub-Kerckhoffs), so that he reads the failure as breakage rather
than as a deliberate move; and it needs him ignorant of NNPP's consequence
(sub-NNPP), so that he does not reason to an undisclosed backup. The second is
not implied by the first, and this is the design's central weakness: \emph{even
total success at the first leaves the encounter open.} ``This device is broken''
does not entail ``this person has no bitcoin''; it entails that the instrument
for reaching it is broken, whose natural sequel is the question \emph{where is
your backup?} --- addressed to a $\Victim$ who is still present and still knows.

So the self-destruct fails against an informed $\Attacker$ for the ordinary
Kerckhoffs reason, and against an ignorant one it fails anyway, because the
belief it manufactures does not discharge the encounter. It is the one obscure
scheme with no branch on which it wins.

We did not have to derive that regress: a vendor documents it. Coldcard ships
trick PINs whose actions include a duress wallet --- ``a personal safety
feature\ldots{} If you must reveal a PIN under duress, give the duress PIN
instead of your Main PIN'' --- alongside \emph{Brick Self} and \emph{Wipe
Seed}~\cite{coldcard-trick-pins}. Its firmware documentation then addresses the
case where the decoy is not believed, and its advice is to escalate to
destruction: ``you may be asked what the actual duress PIN is, while under
duress. We suggest providing the `brickme' PIN in that
case''~\cite{coldcard-pin-entry}. That is the manufacturer contemplating an
$\Attacker$ who knows decoys exist, and recommending, as the answer, an
irreversible erasure performed with $\Victim$ still in the room and still able
to be asked for a backup --- the ignorance-count argument above, arrived at
independently and published as guidance. The same document instructs the holder
to ``triple-check your backups are correct'' before enabling that PIN, so the
backup presupposition of \S\ref{sub:selfdestruct} is not peculiar to one vendor
either.
\end{remark}

\begin{remark}[The class, and what each vendor actually claims for it]
\label{rem:vendor-claims}
Three vendors ship the pair, which is what makes this a design consensus rather
than one firm's misstep; but they do not make the same claims, and the
difference is worth recording precisely. Blockstream markets the \emph{erasure}
for physical threat while claiming nothing of the kind for its passphrase
wallets. Coldcard markets \emph{both} halves, describing the duress wallet as a
safety feature and the brick PIN as the fallback when it fails. Trezor is the
mirror image of Blockstream: its passphrase wallets are sold on deniability,
while its wipe code --- which the vendor itself calls a ``self-destruct'' PIN,
and which erases the recovery seed with it --- is documented without any mention
of duress, coercion or physical threat~\cite{trezor-wipe-code}. That vendor was
moreover asked, in January 2022, to make the wipe \emph{concealable}: a
user-defined message on entry, so that an erasure could be presented as a
hardware fault rather than announcing itself. The request was closed as \emph{not
planned}, and the device still reports the wipe
plainly~\cite{trezor-wipe-deniability}. Component~\textbf{(b)} is therefore
not intrinsic to the self-destruct class at all: it is a presentation one vendor
builds and another, asked for it directly, refused. What \S\ref{sub:selfdestruct}
analyses under~\textbf{(b)} is specific to the ``Internal Error'' design, and
a wipe that announces itself carries~\textbf{(a)} alone --- which, by the
argument above, is no coercion resistance at all, and is not sold as any. Two of the three
therefore decline to make the coercion claim on one of their two mechanisms. We
note it because the criterion is not advanced against vendors, and a vendor that
ships a mechanism without claiming it defeats coercion has not made the error
this section is about.
\end{remark}

\begin{remark}[Two bluffs the erasure invites]\label{rem:erase-bluffs}
An $\Attacker$ who knows the design (Assumption~\ref{ass:power}) need not walk
into the erasure at all. Two alternatives stand open to him, each a wager rather
than a certainty:
\begin{enumerate}[nosep,leftmargin=1.6em]
  \item \emph{Take the device hostage up front.} Rather than enter any PIN,
        seize the device and hold it as a hostage token (\S\ref{sub:hostage},
        Remark~\ref{rem:hostage}). An erase PIN is irrelevant to a coercer who
        never authenticates; and by Remark~\ref{rem:hostage-probe} the terms on
        which $\Victim$ ransoms the device are themselves a probe for the very
        backups $\Victim$ cannot disprove.
  \item \emph{Coerce against deployment.} Brutalize $\Victim$ to deter the PIN's
        use in the first place, wagering that a holder under sufficient duress
        will not spend the one irreversible move available to them.
\end{enumerate}
Neither wager is sure to pay, and that is precisely the point: both are bluff
games, and the holder who has bet on an erasure has entered one without holding
the means to end it. A $\Victim$ who erases and then \emph{truthfully} denies
retaining any backup arrives at the position the registry already records --- the
streamer whose denial was simply disbelieved, and who was tortured for hours
regardless~\cite{ratirl} (\S\ref{sub:empirical}).
\end{remark}

The ``Internal Error'' disguise inherits the same objection from a different
direction: it is a bet on $\Attacker$'s ignorance of a consumer product, and it
degrades exactly as that product's own public documentation becomes known
(\S\ref{sub:fourprops}). The class therefore lands in the \textbf{naive} tier of
Definition~\ref{def:timegate} precisely as that definition describes it --- an
obscure self-destruct, safe only against a sub-Kerckhoffs attacker --- and it is
worth keeping the self-destruct distinct from the \emph{cancel} key named
alongside it there. A cancel reverses a pending action and thereby names a second
racer; a self-destruct destroys a path and names none. Only the former is subject
to Lemma~\ref{lem:drl}, which is why the two fail differently despite sharing a
tier. Remark~\ref{rem:po-variants} has already disposed of the tier as a
whole, by the dilemma that an $\Attacker$ whose ignorance can be assumed makes
the elaborate mechanism unnecessary, and one whose ignorance cannot makes it
ineffective; the present subsection is that dilemma instantiated in a shipping
product.

\medskip\noindent\textbf{The companion decoy, classified separately.} The same
device ships a second and quite different obscure mechanism, and the two should
not be run together. A BIP-39 passphrase ``is added to your Jade's
recovery phrase or SeedQR and is used to create your wallet''; ``any number of
passphrases can be entered to change the wallet you will log into''; and the
device ``will not store your passphrase''~\cite{jade-passphrase}. Distinct
passphrases therefore yield distinct wallets from one recovery phrase, and the
familiar community practice is to keep a low-balance wallet at the empty
passphrase to be surrendered under coercion.

Two observations. First, and to the vendor's credit, \emph{Blockstream does not
market it this way}: the passphrase documentation claims no deniability, decoy
or duress use, in pointed contrast to the erase-PIN pages, which name physical
threat explicitly. The duress framing here is the community's,
not the manufacturer's, and the distinction is worth preserving in any account of
who claimed what.

Second, the mechanism is the plausible-deniability layer already treated in
\S\ref{sub:fourprops} and dismissed for the same Kerckhoffs reason: it rests on
$\Attacker$'s ignorance, and an $\Attacker$ who knows that passphrase wallets
exist simply persists past the surrendered one, with NNPP guaranteeing $\Victim$
cannot close the question. It is not, however, the \emph{same} failure as the
self-destruct's, and by Remark~\ref{rem:two-ignorances} it is the less bad of the
two: the decoy at least has a branch on which it wins, since an $\Attacker$ who
believes it has been paid and can leave. The self-destruct has none. That a
single product ships both, marketing only the weaker one for physical threat, is
the clearest illustration in this section of how little the design space has been
mapped against the Criterion.

\begin{remark}[The decoy is sold, and coached for credibility]\label{rem:decoy-market}
The decoy is not merely a practice the community tolerates; it is a product and a
curriculum. A Spanish-language self-custody consultancy lists, among paid service
tiers, ``\emph{Passphrase (tu b\'oveda se\~nuelo)} --- Billetera se\~nuelo +
billetera real. Protecci\'on ante coacci\'on/coerci\'on'' --- a decoy wallet
marketed, in as many words, as protection against coercion~\cite{tlt-servicios}.
Published guidance goes further and coaches the performance: one guide has the
holder keep a small balance on the bare seed as ``what you'd hand over under
physical coercion'', and adds that ``for this to work, the decoy balance must be
\emph{believable} --- enough to look like a real wallet, not enough to devastate
you if lost''~\cite{bsg-passphrase}.\ifanon\else{} The counsel is also delivered
live: one of us appeared as guest on a self-custody
programme~\cite{ong-autocustodia} in which the host worked through a
decoy-construction walkthrough over the guest's contrary advice.\fi

That last instruction is the naive tier stated as a training objective. A decoy
need be believable only to an $\Attacker$ who might disbelieve it, and the
counsel is silent on what happens when he does --- which, by
Definition~\ref{def:naive-cr}, is the whole question. A better-rehearsed lie
raises the chance $\Attacker$ takes the bait at the margin; it does not create a
branch on which the scheme wins against an $\Attacker$ who knows decoys exist,
and NNPP (Assumption~\ref{lem:nnpl}) guarantees $\Victim$ cannot close that
question however well the performance goes. What is sold as a mechanism is a
wager on the attacker's ignorance, priced and taught as though it were a control.
\end{remark}

\begin{remark}[The Obscurity Paradox --- non-obscure designs strengthen obscure ones]
\label{rem:obscurity-paradox}
TGPO enters neither bluff. There is nothing to erase, the deep secret having
never been on the device (Theorem~\ref{thm:floor}), and so nothing whose
non-deployment can be coerced. But the comparison yields more than the absence
of a failure mode, and the surplus runs in a direction the obscurity argument
would not predict.

An obscure defense asks $\Attacker$ to believe a \emph{claim} --- that this is
the only wallet, that nothing further remains --- against both an incentive to
disbelieve it and, by NNPP, no means of settling it. A non-obscure defense
changes what making that claim costs, and it does so by the mechanism set out in
Remark~\ref{rem:adversary-education}: an $\Attacker$ who credits that an
effective mechanism is in place is crediting that lying is not $\Victim$'s only
available tool, and a holder with effective alternatives has correspondingly
less reason to resort to one.

The ordering this inverts is worth stating plainly, and we name it \textbf{the
Obscurity Paradox}: \emph{a publicly specified, Kerckhoffs-clean design improves
the standing of precisely those obscure denials it was built to make
unnecessary.} Obscurity remains unsound as a primary defense, for the reasons
\S\ref{sub:fourprops} gives. It is not, however, indifferent which design it
accompanies --- a denial made from within a mechanism that removes the need to
lie is a stronger denial than the same words spoken from within one that
requires it.
\end{remark}

\subsection{Registry method and coding tables}
\label{app:classify-registry}

\noindent \S\ref{sub:empirical} reports the headline findings of the coding. The
method, the full tables and the caveats are these.

\emph{Short attacks dominate; the registry over-counts the dramatic.} Most wrench
attacks are \emph{hand-over-now} events --- a threat, an immediate transfer, minutes
on scene; the long, cinematic kidnapping is the exception. A media-sourced registry
inverts this, because a mutilation-kidnapping is news where a knifepoint mugging is
not, so the kidnap share coded below \emph{over}-states the long tail rather than
establishing it. That TKBA is realistically \emph{less} effective against a
long-duration attack --- where hours of pen-and-paper effort under coercion
\emph{might} walk a deep stage --- is granted, but it is \emph{speculative} pending
the field study of \S\ref{sec:limits}. What is \emph{not} speculative is that the
fielded alternatives fail; the cases below make that concrete.

\begin{center}\small
\begin{tabular}{@{}lrr@{}}
\toprule
\textbf{Signal (headline keyword coding)} & \textbf{Incidents} & \textbf{Share} \\
\midrule
Killed / murdered victim \emph{(a lower bound)} & $18$  & $5.1\%$  \\
Kidnap / hostage / ransom (long-hold)           & $139$ & $39.5\%$ \\
Armed robbery at gunpoint (short-hold)          & $101$ & $28.7\%$ \\
Home invasion / break-in / raid                 & $40$  & $11.4\%$ \\
Torture or physical violence                    & $82$  & $23.3\%$ \\
\bottomrule
\end{tabular}
\end{center}

\noindent Categories overlap (a kidnapping may also be violent or fatal) and do not
sum to $100\%$; shares are of all $352$ incidents, and the modality labels are a
headline proxy, not a measured duration. The sample is media-reported and
survivorship-biased --- deterred, unreported, and silently-fatal cases are absent ---
so every figure is directional and the fatality share in particular is a floor.

\medskip\noindent\textbf{The mix is not universal, and the Lemma's bite tracks it.} The
global shares above average over jurisdictions whose compositions differ by an order of magnitude.
Taking the ratio of long-hold to short-hold codings, $K{:}G$:

\begin{center}\small
\begin{tabular}{@{}lrrrr@{}}
\toprule
\textbf{Jurisdiction} & \textbf{$n$} & \textbf{Kidnap} & \textbf{Gunpoint} & \textbf{$K{:}G$} \\
\midrule
Brazil \emph{(small $n$)}  & $12$  & $75.0\%$ & $8.3\%$  & $9.0$ \\
France                     & $61$  & $57.4\%$ & $8.2\%$  & $7.0$ \\
\emph{All incidents}       & $352$ & $39.5\%$ & $28.7\%$ & $1.4$ \\
United States              & $59$  & $28.8\%$ & $40.7\%$ & $0.7$ \\
Russia \emph{(small $n$)}  & $10$  & $30.0\%$ & $50.0\%$ & $0.6$ \\
\bottomrule
\end{tabular}
\end{center}

\noindent The comparison that carries weight is France against the United States: the two largest
national subsets, of near-identical size ($61$ and $59$), with \emph{inverted} compositions --- a
tenfold difference in $K{:}G$. This is the sharpest available check on the media-inflation caveat
just stated. That caveat explains a \emph{level} --- why the registry over-counts dramatic incidents
everywhere --- but it cannot explain this \emph{ordering} without the further claim that French
reporting is ten times more kidnap-selective than American, in a market that is at least as heavily
covered. The parsimonious reading is that modality composition genuinely varies with local
conditions --- firearm availability, the presence of organized kidnap-for-ransom capability --- as
it does for kidnapping outside this domain entirely.

Two confounds remain and are not dismissed. Outside France and the United States the subsets are
too small to carry an argument, so the other rows are directional only. And the coding is applied to
English-language headlines, several of them translations, so a systematic vocabulary difference ---
French sources reaching for ``kidnapped'' where American ones reach for ``robbed at gunpoint'' ---
would produce part of this spread without any difference in the underlying events.

If the composition is real, it matters for scope rather than for validity. The Deadly-Race Lemma
(Lemma~\ref{lem:drl}) is a statement about the \emph{released} holder, and release is the trigger
(Remark~\ref{rem:release-trigger}); a hand-over-now robbery at gunpoint approximates the clean
$p_A = 1$ endpoint --- the attacker takes what is reachable and leaves --- whereas a multi-hour hold
ending in a deliberate release is exactly the gray outcome. Coercion risk is therefore not merely
higher in some jurisdictions but \emph{differently shaped}, and the No-Gray-Area Criterion
(Criterion~\ref{prin:nga}) does its work where long holds dominate. That is the regime the record is
currently growing into: France alone accounts for a majority of the incidents recorded in 2026.

\subsection{The proxy-victim finding}
\label{app:classify-proxy}

\medskip\noindent\textbf{The attacker often does not need the holder at all.} A fourth
signal falls outside the keyword table because it was coded by hand: in $12$ of the $352$
incidents ($3.4\%$) the person seized was \emph{not} a holder but a relative held to compel
one --- a parent, a spouse, an adult child.\footnote{Keyword screening surfaced $16$ candidates;
each was read individually and four were discarded as false positives, including one in which
``son of'' named a \emph{perpetrator} rather than a victim. \texttt{registry-analysis/} carries the
screen and the full audited list. Because the coding reads a one-line headline, an incident that
does not name a relative is no evidence that none was involved: $3.4\%$ is a floor, and the coding
can only understate it.} Two features of this subset matter more than its size.

It is \emph{recent and concentrating}: eleven of the twelve fall in 2023 or later, and \emph{nine}
occur in a single jurisdiction (France) whose recorded incidents rose from about one a year through
2024 to $22$ in 2025 and $34$ by August 2026. Whatever share of that rise is genuine incidence
rather than reporting intensity, the tactic is not a historical scatter; it is being adopted.

It is \emph{not} the hostage token of \S\ref{sub:hostage}, and the difference matters. There the
leverage is a \emph{seizable object} the protocol requires, which devicelessness removes by having
no such object; here it is a \emph{person}, whom no custody design renders unseizable. What this
subset shows in the field is clause~\ref{d:nodelegatecoerce} failing: such an attacker needs no
cryptographic access and no cooperation a protocol could withhold --- only leverage, and a relative
supplies it. This bears directly on the delegated route~(\ref{r:remotewinner}), whose
clearance rests on the deciding party lying beyond $\Attacker$'s coercive reach. The delegate a
holder is most likely to actually appoint --- a spouse, a parent, an adult child --- is precisely
the party this subset shows being seized. The route is not thereby refuted: it is bounded. A
delegate who is reachable is not a second racer but a second victim, so \ref{r:remotewinner}
clears the bar only for delegates whose unreachability is structural rather than assumed --- and
that requirement, not the arrangement's convenience, is what the data prices.

\subsection{Two cases in point}
\label{app:classify-cases}

\medskip\noindent\textbf{A case in point --- Balland.} The January 2025 abduction of
David Balland, a co-founder of the hardware-wallet maker Ledger~\cite{balland}, is
instructive precisely because it is an \emph{embarrassment to the self-custody creed}.
A professional champion of self-custody was seized from his home and mutilated --- a
finger severed to force payment --- and neither wealth nor prominence prevented it. He
was freed by an elite \emph{state} police unit, and the ransom --- paid not by Balland
but by an associate --- was clawed back largely because part of it moved in a
\emph{centralized} stablecoin whose issuer could freeze it. A self-custody pioneer
made whole by the state and by centralization: two mechanisms external to, and in
tension with, the creed. Only the issuer freeze is independently verifiable; the
remainder rests on the parties' accounts, according to which roughly $5\%$ was never recovered.
The case also indicts the delegated route (\ref{r:remotewinner}): an associate
\emph{did} pay, which is exactly the compliance clause~\ref{d:nodelegatecoerce}
assumes away --- a real instance of the indirect coercion,
routed through the victim's body to a third party who will not refuse. For a brief
hand-over-now robbery a remote delegate is unreachable in time and so seems to help;
but merely \emph{not carrying} the device --- leaving the wallet in a home vault ---
does the same, while for the long attack the delegate is not out of reach but is the
ransom channel itself. The lesson is not that Balland erred: it is that, against a
determined long-hold attacker, the fielded options collapse to state rescue and
centralized freezing, whereas the tacit route removes the leverage --- no standing
device, no hostage token, no delegate to pay --- rather than relying on ex-post
recovery.

\medskip\noindent\textbf{A case in point --- RATIRL.} The registry coded above records
the obscurity route's failure directly. In one of its 2026 incidents (Sweden, 13
January)~\cite{ratirl} a League of Legends streamer was bound and tortured for hours by
armed intruders demanding cryptocurrency he denied holding; unable to credit the
denial, they persisted long past any value the robbery could rationally return,
leaving finally with only his Pok\'emon-card collection. This is
Corollary~\ref{cor:disclosure} in the field: where the community norm is ``deny that
you own Bitcoin,'' a coerced denial is \emph{pre-discredited} --- presumed to be the
denial-script --- so a Kerckhoffs-educated $\Attacker$ gets no exculpatory update, and
the rational continuation of an interrogation that opens with a denial is \emph{more}
coercion, not less. The case fixes the
\emph{sign}: obscurity does not merely fail to protect --- it can convert a robbery
into a torture --- \emph{negative} coercion resistance, not a mere absence. And it does so on \emph{every}
reading of the denial, which forks three ways --- each indicting obscurity.
\textbf{(i)}~If the denial was a defensive lie (there was a stash), secrecy still did
not avert the assault --- a textbook Kerckhoffs failure --- and his only residual
recourse is to \emph{double down} on it and wager against a repeat. \textbf{(ii)}~If it
was true (no stash), he bore obscurity's \emph{negative externality}: a genuine
non-holder attacked only because real holders are coached to look exactly like him ---
the very coaching that collapses the statistical holder/non-holder distinction.
\textbf{(iii)}~If the truth is in between --- a real but smaller-than-presumed stash,
too small to be worth the judicial exposure the attack incurred --- \emph{both}
pathologies operate at once: others' obscurity redirected the attack onto him, and his
own failed to stop it. That all three are simultaneously plausible is itself the
verdict: in the aggregate, across a growing population of holders and lookalikes, each
is already occurring --- and the obscurity culture, exactly as Kerckhoffs predicts,
prevents none of them.
